\documentclass[pdflatex,sn-mathphys-num]{sn-jnl}


\usepackage{graphicx}
\usepackage{multirow}%
\usepackage{amsmath,amssymb,amsfonts}%
\usepackage{amsthm}%
\usepackage[title]{appendix}%
\usepackage{xcolor}%
\usepackage{textcomp}%
\usepackage{manyfoot}%
\usepackage{booktabs,bookmark}
\usepackage{algorithm}%
\usepackage{algorithmicx}%
\usepackage{algpseudocode}%
\usepackage{listings}%
\usepackage{microtype}
\usepackage{csquotes}
\usepackage{enumitem}
\usepackage{hyperref}
\usepackage{cleveref}
\usepackage{siunitx}
\usepackage{derivative}
\usepackage{braket}

\renewcommand{\arraystretch}{1.2}
\sisetup{group-digits=true,
  group-separator={\,},
  separate-uncertainty
}

\hypersetup{
  colorlinks=true,
  linkcolor=black,
  urlcolor=blue,
  pdftitle={Identical Particles in Infinite Well},
  pdfauthor={Avila, Herrera, Rodríguez},
  pdfsubject={Quantum Mechanics II},
  pdfkeywords={Homework 7}
}
\urlstyle{same}

\theoremstyle{thmstyleone}%
\newtheorem{theorem}{Theorem}%
\newtheorem{proposition}[theorem]{Proposition}%

\theoremstyle{thmstyletwo}%
\newtheorem{example}{Example}%
\newtheorem{remark}{Remark}%

\theoremstyle{thmstylethree}%
\newtheorem{definition}{Definition}%

\raggedbottom

\begin{document}

\title[Angular Momentum Representation For Free perticle ]{\textbf{Angular Momentum Representation For Free perticle}}
\author[]{\fnm{Julian} \sur{Avila}}
\author[]{\fnm{Laura} \sur{Herrera}}
\author[]{\fnm{Sebastian} \sur{Rodríguez}}
\affil[]{\orgname{Universidad Distrital Francisco José de Caldas}}

\abstract{
n
}
\keywords{
  a
}

\maketitle

\section{Problem Setting}

We seek a complete formulation of the angular-momentum representation for a
non-relativistic free particle. The central objective is the diagonalization of
the Hamiltonian
\begin{equation}
  \hat{H} = \frac{\hat{\mathbf{p}}^{\,2}}{2m}
\end{equation}
within the basis defined by the Complete Set of Commuting Observables (CSCO):
$\{ \hat{H}, \hat{L}^2, \hat{L}_z \}$. This corresponds to solving the
time-dependent Schr\"odinger equation (TDSE) by exploiting the isotropy of free
space.

The formulation is motivated by three theoretical necessities:
\begin{enumerate}
  \item \textbf{Symmetry Realization:} To represent the rotation group
    $\mathrm{SO}(3)$
  via its generators acting on the Hilbert space of the free particle.
  \item \textbf{Conservation Laws:} To explicitly link the invariance of
  $\hat{H}$ under spatial rotations to the conservation of angular momentum.
  \item \textbf{Separation of Variables:} To derive the natural basis for
  spherical geometries, generalizing to central potential problems.
\end{enumerate}

\subsection{Physical Scenario}

Consider a spinless particle of mass $m$ in Euclidean space $\mathbb{R}^3$. The
absence of external potentials implies the Hamiltonian is purely kinetic. The
isotropy of the system ensures that the generator of rotations, the angular
momentum vector operator $\hat{\mathbf{L}}$, commutes with the Hamiltonian:
\begin{equation}
  [\hat{H}, \hat{\mathbf{L}}] = 0.
\end{equation}
Consequently, stationary states may be labeled simultaneously by energy ($E_k$),
total angular momentum ($\ell$), and its projection ($m$), i.e., $\hat{H}$,
$\hat{L}^2$ and $\hat{L}_z$ admit simultaneous eigenstates.

\subsection{Construction Objectives}

The project proceeds through the following formal steps:
\begin{itemize}
  \item \textbf{Algebraic Definition:} Derivation of the operator
  $\hat{\mathbf{L}} = \hat{\mathbf{r}} \times \hat{\mathbf{p}}$ utilizing
  fundamental canonical commutation relations (CCR).
  \item \textbf{Lie Algebra $\mathfrak{so}(3)$:} Verification of the closure
  relation $[\hat{L}_i, \hat{L}_j] = i\hbar\, \epsilon_{ijk}\hat{L}_k$,
  establishing the connection to the structure constants of the rotation group.
  \item \textbf{Spectral Analysis:} Determination of the spectrum for
  $\hat{L}^2$ and $\hat{L}_z$, yielding the spherical harmonics
  $Y_{\ell m}(\theta, \phi)$ as coordinate representations of the angular
  eigenbasis.
  \item \textbf{Basis Transformation:} Construction of the unitary
  transformation connecting the linear momentum basis $\ket{\mathbf{k}}$ to the
  spherical basis $\ket{k,\ell,m}$.
\end{itemize}

\subsection{Mathematical Prerequisites}

The derivation assumes fluency in:
\begin{itemize}
  \item \textbf{Abstract Formalism:} Hilbert space $\mathcal{H}$, Dirac
  notation, and spectral theory of Hermitian operators.
  \item \textbf{Symplectic Structure:} Canonical commutators
  $[\hat{x}_i, \hat{p}_j]=i\hbar\,\delta_{ij}$.
  \item \textbf{Integral Transforms:} Fourier analysis relating position
  ($\mathbf{r}$) and momentum ($\mathbf{p}$) representations.
  \item \textbf{Special Functions:} Properties of Spherical Harmonics and
  Spherical Bessel functions.
\end{itemize}

\subsection{Formal Outcome}

The construction culminates in the general solution to the TDSE, expressed as a
superposition of spherical waves:
\begin{equation}
  \ket{\psi(t)} = \sum_{\ell=0}^{\infty} \sum_{m=-\ell}^{\ell} \int_0^\infty
  \mathrm{d}k\, k^2\, \psi_{k\ell m}(0)\, e^{-i \frac{\hbar k^2}{2m} t}
  \ket{k, \ell, m}.
\end{equation}
This expansion connects the abstract eigenkets to the coordinate representation
via the spherical Bessel transform.


\section{Starting Point: Fundamental Formalism}

We formally define the Hilbert space $\mathcal{H} = L^2(\mathbb{R}^3)$ and the
dynamical structure of the free particle. The following definitions establish
the algebraic and spectral foundations required for the angular-momentum
decomposition.

\subsection{Dynamics and Unitary Evolution}

The dynamics are governed by the self-adjoint Hamiltonian operator $\hat{H} =
\frac{\hat{\mathbf{p}}^2}{2m}$. The state vector $\ket{\psi(t)}$ evolves
deterministically from an initial state $\ket{\psi_0}$ via the strongly
continuous unitary group $\hat{U}(t)$:
\begin{equation}
  \ket{\psi(t)} = \hat{U}(t)\ket{\psi_0} =
  \exp\left(-\frac{i}{\hbar}\hat{H}t\right)\ket{\psi_0}.
\end{equation}
This abstract evolution implies the Schrödinger equation $i\hbar \partial_t
\ket{\psi} = \hat{H}\ket{\psi}$ via Stone's Theorem.

\subsection{Continuous Representations}

We utilize two continuous bases, position $\{\ket{\mathbf{r}}\}$ and linear
momentum $\{\ket{\mathbf{p}}\}$, normalized to the Dirac distribution.

\subsubsection{Basis Definitions and Orthogonality}
The basis vectors satisfy the eigenvalue equations and completeness relations
(Resolution of Identity $\hat{\mathbb{I}}$):
\begin{align}
  \hat{\mathbf{r}}\ket{\mathbf{r}} &= \mathbf{r}\ket{\mathbf{r}}, \quad &
  \int_{\mathbb{R}^3} \mathrm{d}^3r \, \ket{\mathbf{r}}\bra{\mathbf{r}} &=
  \hat{\mathbb{I}}, \\
  \hat{\mathbf{p}}\ket{\mathbf{p}} &= \mathbf{p}\ket{\mathbf{p}}, \quad &
  \int_{\mathbb{R}^3} \mathrm{d}^3p \, \ket{\mathbf{p}}\bra{\mathbf{p}} &=
  \hat{\mathbb{I}}.
\end{align}
The inner product defines the transformation kernel (Fourier kernel) between the
canonically conjugate variables:
\begin{equation}
  \braket{\mathbf{r}}{\mathbf{p}} = \frac{1}{(2\pi\hbar)^{3/2}} \exp\left(
  \frac{i}{\hbar} \mathbf{p}\cdot\mathbf{r} \right).
\end{equation}

\subsubsection{Wavefunction Representations}
The projections of the state vector yield the coordinate and momentum space
wavefunctions:
\begin{itemize}
  \item \textbf{Coordinate Space:} $\psi(\mathbf{r}, t) \equiv
    \braket{\mathbf{r}}{\psi(t)}$. The momentum operator acts as the generator
    of spatial translations:
    \begin{equation}
      \bra{\mathbf{r}}\hat{\mathbf{p}}\ket{\psi} = -i\hbar \nabla
      \psi(\mathbf{r}).
    \end{equation}
  \item \textbf{Momentum Space:} $\tilde{\psi}(\mathbf{p}, t) \equiv
    \braket{\mathbf{p}}{\psi(t)}$. Since $\hat{H}$ is diagonal in this basis,
    the time evolution is strictly a phase modulation:
    \begin{equation}
      \tilde{\psi}(\mathbf{p}, t) = \exp\left( -\frac{i}{\hbar}
      \frac{p^2}{2m} t \right) \tilde{\psi}(\mathbf{p}, 0).
    \end{equation}
\end{itemize}

\subsection{Spectral Decomposition and Density of States}

For the free particle, the energy $E$ is degenerate. The transition from
momentum integration to energy integration requires the spectral density
function. Using the dispersion relation $E_p = p^2/2m$:
\begin{equation}
  \int \mathrm{d}^3p = \int_0^\infty p^2 \mathrm{d}p \int_{S^2}
  \mathrm{d}\Omega.
\end{equation}
Changing variables from $p$ to $E$ introduces the density of states factor via
the identity:
\begin{equation}
  \mathrm{d}p = \frac{\mathrm{d}p}{\mathrm{d}E} \mathrm{d}E = \frac{m}{p}
  \mathrm{d}E = \sqrt{\frac{m}{2E}} \mathrm{d}E.
\end{equation}
This transformation is critical for isolating the radial (energy) dependence
from the angular dependence in the final construction.

\subsection{Symmetry Generators}

The continuous symmetries of the system are generated by Hermitian operators,
implying specific conservation laws via Ehrenfest's Theorem,
$\frac{\mathrm{d}}{\mathrm{d}t}\braket{\hat{A}} = \frac{1}{i\hbar}\braket{[\hat{A},
\hat{H}]}$.

\begin{enumerate}
  \item \textbf{Spatial Translations:} Generated by $\hat{\mathbf{p}}$. Since
    $[\hat{\mathbf{p}}, \hat{H}] = 0$, linear momentum is conserved.
    \begin{equation}
      \hat{T}(\mathbf{a}) = \exp\left( -\frac{i}{\hbar} \mathbf{a} \cdot
      \hat{\mathbf{p}} \right).
    \end{equation}
  \item \textbf{Galilean Boosts:} Generated by $\hat{\mathbf{r}}$ (up to mass
    scaling). This shifts the momentum spectrum:
    \begin{equation}
      \hat{T}_{\mathbf{p}}(\mathbf{q}) = \exp\left( \frac{i}{\hbar}
      \mathbf{q} \cdot \hat{\mathbf{r}} \right).
    \end{equation}
\end{enumerate}


\section{The Angular Momentum Operator}

In classical mechanics, the angular momentum of a particle with position vector
$\mathbf{r}$ and linear momentum $\mathbf{p}$ is defined as the vector product
$\mathbf{L}_{cl} = \mathbf{r} \times \mathbf{p}$. To transition to quantum
mechanics, we apply the correspondence principle, promoting the dynamical
variables to linear operators acting on the Hilbert space: $\mathbf{r} \to
\hat{\mathbf{r}}$ and $\mathbf{p} \to \hat{\mathbf{p}}$.

However, since $\hat{\mathbf{r}}$ and $\hat{\mathbf{p}}$ are non-commuting
operators, the definition of the vector product must be handled with algebraic
rigor to ensure the resulting operator is Hermitian (measurable) and
unambiguous.

\subsection{Standard Vector Definition}

In 3D Euclidean space, we define the operator $\hat{\mathbf{L}}$ as the vector
product of the position and momentum operators:
\begin{equation}
  \hat{\mathbf{L}} \equiv \hat{\mathbf{r}} \times \hat{\mathbf{p}} = -i\hbar
  (\mathbf{r} \times \nabla).
\end{equation}
Using the Levi-Civita pseudotensor $\epsilon_{ijk}$, the components are defined
as:
\begin{equation} \label{eq:L_tensor_def}
  \hat{L}_i = \epsilon_{ijk} \hat{x}_j \hat{p}_k,
\end{equation}
where Einstein summation is implied. This definition ensures that
$\hat{\mathbf{L}}$ acts as a vector operator under spatial rotations.

\subsection{Geometric Algebra Perspective (Clifford Algebra)}

To provide a coordinate-independent definition valid in any dimension, we employ
the Geometric Product within the algebra $\mathcal{C}\ell_{3,0}(\mathbb{R})$.
The product of the vector operators $\hat{\mathbf{r}}$ and $\hat{\mathbf{p}}$
decomposes into a scalar (symmetric) and bivector (antisymmetric) part:
\begin{equation}
  \hat{\mathbf{r}}\hat{\mathbf{p}} = \hat{\mathbf{r}} \cdot \hat{\mathbf{p}} +
  \hat{\mathbf{r}} \wedge \hat{\mathbf{p}}.
\end{equation}
The physically significant quantity is the bivector $\hat{\mathcal{L}}$,
representing the oriented plane of rotation:
\begin{equation}
  \hat{\mathcal{L}} \equiv \hat{\mathbf{r}} \wedge \hat{\mathbf{p}} =
  \frac{1}{2}(\hat{\mathbf{r}}\hat{\mathbf{p}} -
  \hat{\mathbf{p}}\hat{\mathbf{r}}).
\end{equation}
The standard vector operator $\hat{\mathbf{L}}$ is recovered as the Hodge
dual of this bivector. Using the unit pseudoscalar $I = \hat{e}_1 \hat{e}_2
\hat{e}_3$:
\begin{equation}
  \hat{\mathbf{L}} = -i I \hat{\mathcal{L}}.
\end{equation}
This formulation reveals that spin and orbital angular momentum share the same
bivector structure, differing only in the vector space they act upon.

\subsection{Hermiticity and Observability}

For $\hat{\mathbf{L}}$ to represent a physical observable, it must be Hermitian.
We examine the adjoint of the component definition (\ref{eq:L_tensor_def}):
\begin{equation}
  \hat{L}_i^\dagger = (\epsilon_{ijk} \hat{x}_j \hat{p}_k)^\dagger.
\end{equation}
Using the property $(\hat{A}\hat{B})^\dagger = \hat{B}^\dagger \hat{A}^\dagger$
and the fact that $\hat{\mathbf{x}}$ and $\hat{\mathbf{p}}$ are Hermitian
($\hat{x}_j^\dagger = \hat{x}_j$, $\hat{p}_k^\dagger = \hat{p}_k$) and
$\epsilon_{ijk}$ is real:
\begin{equation}
  \hat{L}_i^\dagger = \epsilon_{ijk} \hat{p}_k^\dagger \hat{x}_j^\dagger =
  \epsilon_{ijk} \hat{p}_k \hat{x}_j.
\end{equation}
The non-commutativity of $\hat{\mathbf{x}}$ and $\hat{\mathbf{p}}$ is resolved
by the canonical commutation relation $[\hat{x}_j, \hat{p}_k] = i\hbar
\delta_{jk}$:
\begin{equation}
  [\hat{x}_j, \hat{p}_k] = i\hbar \delta_{jk} \implies \hat{p}_k \hat{x}_j =
  \hat{x}_j \hat{p}_k - i\hbar \delta_{jk}.
\end{equation}
Substituting this back into the expression for $\hat{L}_i^\dagger$:
\begin{align}
  \hat{L}_i^\dagger &= \epsilon_{ijk} (\hat{x}_j \hat{p}_k - i\hbar \delta_{jk})
  \nonumber \\
                    &= \epsilon_{ijk} \hat{x}_j \hat{p}_k - i\hbar \epsilon_{ijk}
  \delta_{jk}.
\end{align}
The second term contains the contraction of the antisymmetric tensor
$\epsilon_{ijk}$ with the symmetric Kronecker delta $\delta_{jk}$. Since
$\epsilon_{ijk}$ vanishes if any two indices are equal (which $\delta_{jk}$
forces), we have:
\begin{equation}
  \epsilon_{ijk} \delta_{jk} = \epsilon_{ijj} = 0.
\end{equation}
Therefore:
\begin{equation}
  \hat{L}_i^\dagger = \epsilon_{ijk} \hat{x}_j \hat{p}_k = \hat{L}_i.
\end{equation}
The operator is inherently Hermitian, and no symmetrization (e.g.,
$\frac{1}{2}(\mathbf{r}\times\mathbf{p} - \mathbf{p}\times\mathbf{r})$) is
required because the cross product couples only commuting orthogonal components
(e.g., $x$ and $p_y$).

\subsection{Explicit Component Form}

Expanding the tensor notation, we obtain the explicit forms used in standard
calculations:

\begin{align}
  \hat{L}_x &= \hat{y}\hat{p}_z - \hat{z}\hat{p}_y, \\
  \hat{L}_y &= \hat{z}\hat{p}_x - \hat{x}\hat{p}_z, \\
  \hat{L}_z &= \hat{x}\hat{p}_y - \hat{y}\hat{p}_x.
\end{align}

These expressions confirm that the quantum angular momentum operator preserves
the functional form of its classical counterpart, provided the order of
operators is maintained as $\mathbf{r} \times \mathbf{p}$.

\subsection{The Lie Algebra \texorpdfstring{$\mathfrak{so}(3)$}{so(3)}}

The fundamental commutation relations for the components of $\hat{\mathbf{L}}$
define the structure of the rotation group. Computing $[\hat{L}_x, \hat{L}_y]$:
\begin{align}
  [\hat{L}_x, \hat{L}_y] &= [y p_z - z p_y, z p_x - x p_z] \nonumber \\
                         &= [y p_z, z p_x] + [-z p_y, -x p_z] \quad
                         (\text{other terms commute}) \nonumber \\
                         &= y p_x [p_z, z] + x p_y [z, p_z] \nonumber \\
                         &= y p_x (-i\hbar) + x p_y (i\hbar) \nonumber \\
                         &= i\hbar (x p_y - y p_x) = i\hbar \hat{L}_z.
\end{align}
Generalizing this result yields the Lie algebra of $\mathfrak{so}(3)$:
\begin{equation}
  [\hat{L}_i, \hat{L}_j] = i\hbar \epsilon_{ijk} \hat{L}_k.
\end{equation}
This closure relation confirms that $\hat{\mathbf{L}}$ generates the group of
spatial rotations $\mathrm{SO}(3)$.


\section{Representation of $\hat{L}$ in Position Space}

In the position representation, the Hilbert space is spanned by the eigenkets $|\mathbf{r}\rangle$. The position operator acts multiplicatively, $\hat{\mathbf{r}} \to \mathbf{r}$, and the linear momentum operator acts as the gradient generator.

\subsection{Representation in Cartesian Coordinates}

\subsubsection{Derivation from Fundamental Operators}
Using the position representation of the momentum operator[cite: 42]:
\begin{equation}
    \hat{\mathbf{p}} = -i\hbar \nabla_\mathbf{r},
\end{equation}
we substitute this into the definition of angular momentum derived in the previous section ($\hat{\mathbf{L}} = \hat{\mathbf{r}} \times \hat{\mathbf{p}}$):
\begin{equation}
    \hat{\mathbf{L}} = \hat{\mathbf{r}} \times (-i\hbar \nabla) = -i\hbar (\mathbf{r} \times \nabla).
\end{equation}

\subsubsection{Explicit Components}
Expanding the cross product in Cartesian coordinates using the determinant mnemonic (or the Levi-Civita definition):
\begin{equation}
    \mathbf{r} \times \nabla = 
    \begin{vmatrix} 
    \hat{\mathbf{e}}_x & \hat{\mathbf{e}}_y & \hat{\mathbf{e}}_z \\ 
    x & y & z \\ 
    \partial_x & \partial_y & \partial_z 
    \end{vmatrix} 
    = \hat{\mathbf{e}}_x (y\partial_z - z\partial_y) - \hat{\mathbf{e}}_y (x\partial_z - z\partial_x) + \hat{\mathbf{e}}_z (x\partial_y - y\partial_x).
\end{equation}
Multiplying by the factor $-i\hbar$, we obtain the explicit components[cite: 45]:
\begin{align}
    \hat{L}_x &= -i\hbar \left( y \frac{\partial}{\partial z} - z \frac{\partial}{\partial y} \right), \\
    \hat{L}_y &= -i\hbar \left( z \frac{\partial}{\partial x} - x \frac{\partial}{\partial z} \right), \\
    \hat{L}_z &= -i\hbar \left( x \frac{\partial}{\partial y} - y \frac{\partial}{\partial x} \right). \label{eq:Lz_cartesian}
\end{align}

\subsection{Representation in Spherical Coordinates}

We now transform these operators to spherical coordinates $(r, \theta, \phi)$ defined by:
\begin{align*}
    x &= r \sin\theta \cos\phi, \\
    y &= r \sin\theta \sin\phi, \\
    z &= r \cos\theta.
\end{align*}

\subsubsection{Derivation of $\hat{L}_z$}
To find the expression for $\hat{L}_z$ in spherical coordinates, we apply the Chain Rule to the partial derivative with respect to $\phi$. 
Let $\psi(x,y,z)$ be a function of position. The derivative with respect to $\phi$ is:
\begin{equation}
    \frac{\partial}{\partial \phi} = \frac{\partial x}{\partial \phi}\frac{\partial}{\partial x} + \frac{\partial y}{\partial \phi}\frac{\partial}{\partial y} + \frac{\partial z}{\partial \phi}\frac{\partial}{\partial z}.
\end{equation}
Calculating the partial derivatives of the coordinates:
\begin{itemize}
    \item $\frac{\partial x}{\partial \phi} = \frac{\partial}{\partial \phi}(r \sin\theta \cos\phi) = -r \sin\theta \sin\phi = -y$.
    \item $\frac{\partial y}{\partial \phi} = \frac{\partial}{\partial \phi}(r \sin\theta \sin\phi) = r \sin\theta \cos\phi = x$.
    \item $\frac{\partial z}{\partial \phi} = 0$ (since $z$ does not depend on $\phi$).
\end{itemize}
Substituting these back into the chain rule expression:
\begin{equation}
    \frac{\partial}{\partial \phi} = (-y)\frac{\partial}{\partial x} + (x)\frac{\partial}{\partial y} = x\frac{\partial}{\partial y} - y\frac{\partial}{\partial x}.
\end{equation}
Comparing this result with Eq. (\ref{eq:Lz_cartesian}), we immediately identify the term in parentheses. Thus:
\begin{equation}
    \hat{L}_z = -i\hbar \left( x \frac{\partial}{\partial y} - y \frac{\partial}{\partial x} \right) = -i\hbar \frac{\partial}{\partial \phi}.
\end{equation}

\subsubsection{Derivation of $\hat{L}^2$}
To derive the squared total angular momentum $\hat{L}^2 = \hat{\mathbf{L}} \cdot \hat{\mathbf{L}}$, we avoid the tedious squaring of Cartesian components and instead use the vector properties of the Gradient in spherical coordinates.

\textbf{Step 1: The Gradient in Spherical Coordinates}
The gradient operator $\nabla$ is expressed in the spherical orthonormal basis $(\hat{\mathbf{e}}_r, \hat{\mathbf{e}}_\theta, \hat{\mathbf{e}}_\phi)$ as:
\begin{equation}
    \nabla = \hat{\mathbf{e}}_r \frac{\partial}{\partial r} + \hat{\mathbf{e}}_\theta \frac{1}{r} \frac{\partial}{\partial \theta} + \hat{\mathbf{e}}_\phi \frac{1}{r \sin\theta} \frac{\partial}{\partial \phi}.
\end{equation}

\textbf{Step 2: The Vector $\hat{\mathbf{L}}$ in Spherical Basis}
Using $\mathbf{r} = r \hat{\mathbf{e}}_r$ and calculating the cross product $\hat{\mathbf{L}} = -i\hbar (\mathbf{r} \times \nabla)$:
\begin{equation}
    \hat{\mathbf{L}} = -i\hbar \left[ r\hat{\mathbf{e}}_r \times \left( \hat{\mathbf{e}}_r \frac{\partial}{\partial r} + \hat{\mathbf{e}}_\theta \frac{1}{r} \frac{\partial}{\partial \theta} + \hat{\mathbf{e}}_\phi \frac{1}{r \sin\theta} \frac{\partial}{\partial \phi} \right) \right].
\end{equation}
Using the orthogonality relations $\hat{\mathbf{e}}_r \times \hat{\mathbf{e}}_r = 0$, $\hat{\mathbf{e}}_r \times \hat{\mathbf{e}}_\theta = \hat{\mathbf{e}}_\phi$, and $\hat{\mathbf{e}}_r \times \hat{\mathbf{e}}_\phi = -\hat{\mathbf{e}}_\theta$:
\begin{align}
    \hat{\mathbf{L}} &= -i\hbar \left[ \left( r \frac{1}{r} \right) (\hat{\mathbf{e}}_r \times \hat{\mathbf{e}}_\theta) \frac{\partial}{\partial \theta} + \left( r \frac{1}{r \sin\theta} \right) (\hat{\mathbf{e}}_r \times \hat{\mathbf{e}}_\phi) \frac{\partial}{\partial \phi} \right] \nonumber \\
    &= -i\hbar \left( \hat{\mathbf{e}}_\phi \frac{\partial}{\partial \theta} - \hat{\mathbf{e}}_\theta \frac{1}{\sin\theta} \frac{\partial}{\partial \phi} \right). \label{eq:L_vector_spherical}
\end{align}

\textbf{Step 3: Calculating $\hat{L}^2 = \hat{\mathbf{L}} \cdot \hat{\mathbf{L}}$}
We apply the vector operator $\hat{\mathbf{L}}$ from Eq. (\ref{eq:L_vector_spherical}) to itself. Extreme care must be taken because the unit vectors $\hat{\mathbf{e}}_\theta$ and $\hat{\mathbf{e}}_\phi$ depend on the coordinates, so derivatives acting on them are non-zero.
\begin{equation}
    \hat{L}^2 = (-i\hbar)^2 \left( \hat{\mathbf{e}}_\phi \frac{\partial}{\partial \theta} - \frac{\hat{\mathbf{e}}_\theta}{\sin\theta} \frac{\partial}{\partial \phi} \right) \cdot \left( \hat{\mathbf{e}}_\phi \frac{\partial}{\partial \theta} - \frac{\hat{\mathbf{e}}_\theta}{\sin\theta} \frac{\partial}{\partial \phi} \right).
\end{equation}
Expanding the dot product (remembering that the first operator acts on everything to its right, including unit vectors):
\begin{equation}
    \frac{\hat{L}^2}{-\hbar^2} = \underbrace{\hat{\mathbf{e}}_\phi \cdot \frac{\partial}{\partial \theta} \left( \hat{\mathbf{e}}_\phi \frac{\partial}{\partial \theta} \right)}_{\text{Term 1}} 
    - \underbrace{\hat{\mathbf{e}}_\phi \cdot \frac{\partial}{\partial \theta} \left( \frac{\hat{\mathbf{e}}_\theta}{\sin\theta} \frac{\partial}{\partial \phi} \right)}_{\text{Term 2}} 
    - \underbrace{\frac{\hat{\mathbf{e}}_\theta}{\sin\theta} \cdot \frac{\partial}{\partial \phi} \left( \hat{\mathbf{e}}_\phi \frac{\partial}{\partial \theta} \right)}_{\text{Term 3}} 
    + \underbrace{\frac{\hat{\mathbf{e}}_\theta}{\sin\theta} \cdot \frac{\partial}{\partial \phi} \left( \frac{\hat{\mathbf{e}}_\theta}{\sin\theta} \frac{\partial}{\partial \phi} \right)}_{\text{Term 4}}.
\end{equation}

\textbf{Required Derivatives of Unit Vectors:}
\begin{itemize}
    \item $\partial_\theta \hat{\mathbf{e}}_\phi = 0$.
    \item $\partial_\theta \hat{\mathbf{e}}_\theta = -\hat{\mathbf{e}}_r$ (Orthogonal to $\hat{\mathbf{e}}_\phi$, so dot product vanishes).
    \item $\partial_\phi \hat{\mathbf{e}}_\phi = -\sin\theta \hat{\mathbf{e}}_r - \cos\theta \hat{\mathbf{e}}_\theta$.
    \item $\partial_\phi \hat{\mathbf{e}}_\theta = \cos\theta \hat{\mathbf{e}}_\phi$.
\end{itemize}

\textbf{Evaluating Terms:}

\textit{Term 1:} Since $\partial_\theta \hat{\mathbf{e}}_\phi = 0$:
\[ \hat{\mathbf{e}}_\phi \cdot \left( \hat{\mathbf{e}}_\phi \frac{\partial^2}{\partial \theta^2} \right) = \frac{\partial^2}{\partial \theta^2}. \]

\textit{Term 2:} Contains $\hat{\mathbf{e}}_\phi \cdot \partial_\theta \hat{\mathbf{e}}_\theta$. Since $\partial_\theta \hat{\mathbf{e}}_\theta = -\hat{\mathbf{e}}_r$, the dot product with $\hat{\mathbf{e}}_\phi$ is 0. Cross terms vanish here.

\textit{Term 3:}
\[ \frac{\hat{\mathbf{e}}_\theta}{\sin\theta} \cdot \left( (\partial_\phi \hat{\mathbf{e}}_\phi) \frac{\partial}{\partial \theta} + \hat{\mathbf{e}}_\phi \frac{\partial^2}{\partial \phi \partial \theta} \right). \]
Using $\partial_\phi \hat{\mathbf{e}}_\phi = -\sin\theta \hat{\mathbf{e}}_r - \cos\theta \hat{\mathbf{e}}_\theta$:
\[ \frac{\hat{\mathbf{e}}_\theta}{\sin\theta} \cdot \left( -\cos\theta \hat{\mathbf{e}}_\theta \frac{\partial}{\partial \theta} \right) = -\frac{\cos\theta}{\sin\theta} \frac{\partial}{\partial \theta} = -\cot\theta \frac{\partial}{\partial \theta}. \]

\textit{Term 4:}
\[ \frac{\hat{\mathbf{e}}_\theta}{\sin\theta} \cdot \frac{\partial}{\partial \phi} \left( \frac{\hat{\mathbf{e}}_\theta}{\sin\theta} \frac{\partial}{\partial \phi} \right). \]
The derivative acts on $\hat{\mathbf{e}}_\theta$ and the function. However, $\partial_\phi \hat{\mathbf{e}}_\theta = \cos\theta \hat{\mathbf{e}}_\phi$, which is orthogonal to the pre-factor $\hat{\mathbf{e}}_\theta$. Thus, only the derivative on the scalar function survives:
\[ \frac{1}{\sin^2\theta} \frac{\partial^2}{\partial \phi^2}. \]

\textbf{Assembly:}
Summing the non-vanishing contributions:
\begin{equation}
    \frac{\hat{L}^2}{-\hbar^2} = \frac{\partial^2}{\partial \theta^2} + \cot\theta \frac{\partial}{\partial \theta} + \frac{1}{\sin^2\theta}\frac{\partial^2}{\partial \phi^2}.
\end{equation}
Recognizing that $\frac{\partial^2}{\partial \theta^2} + \cot\theta \frac{\partial}{\partial \theta} = \frac{1}{\sin\theta}\frac{\partial}{\partial \theta}\left(\sin\theta \frac{\partial}{\partial \theta}\right)$, we arrive at the final expression[cite: 48]:

\begin{equation}
    \hat{L}^2 = -\hbar^2 \left[ \frac{1}{\sin\theta}\frac{\partial}{\partial \theta}\left(\sin\theta \frac{\partial}{\partial \theta}\right) + \frac{1}{\sin^2\theta}\frac{\partial^2}{\partial \phi^2} \right].
\end{equation}

\section{Angular Momentum as the Generator of Rotations}

In this section, we establish the fundamental role of the angular momentum operator $\hat{\mathbf{L}}$ as the generator of spatial rotations. We proceed from the definition of the unitary rotation operator to the conservation laws derived from the symmetry of the Hamiltonian.

\subsection{Finite Rotation Operator}

\subsubsection{Definition}
We define the rotation of a physical system by an angle $|\boldsymbol{\theta}|$ around an axis defined by the unit vector $\mathbf{n} = \boldsymbol{\theta}/|\boldsymbol{\theta}|$ using the exponential map. The unitary operator for a finite rotation is defined as[cite: 53]:
\begin{equation}
    \hat{R}(\boldsymbol{\theta}) = \exp\left(-\frac{i}{\hbar} \boldsymbol{\theta} \cdot \hat{\mathbf{L}}\right).
\end{equation}
Since $\hat{\mathbf{L}}$ is Hermitian, $\hat{R}(\boldsymbol{\theta})$ is unitary ($\hat{R}^\dagger = \hat{R}^{-1}$), ensuring the conservation of probability (norm of the state) during rotation.

\subsubsection{Action on Position Kets}
To demonstrate the action of $\hat{R}(\boldsymbol{\theta})$ on the position eigenkets $|\mathbf{r}\rangle$, we consider the transformation of the position operator $\hat{\mathbf{r}}$. Using the Baker-Campbell-Hausdorff lemma or the infinitesimal properties derived below, one can show that:
\begin{equation}
    \hat{R}^\dagger(\boldsymbol{\theta}) \hat{\mathbf{r}} \hat{R}(\boldsymbol{\theta}) = \mathcal{R}(\boldsymbol{\theta}) \hat{\mathbf{r}},
\end{equation}
where $\mathcal{R}(\boldsymbol{\theta})$ is the classical $3\times3$ rotation matrix.
Applying this to a position eigenket $|\mathbf{r}\rangle$:
\begin{equation}
    \hat{\mathbf{r}} \left( \hat{R}(\boldsymbol{\theta}) |\mathbf{r}\rangle \right) = \hat{R}(\boldsymbol{\theta}) \left( \hat{R}^\dagger \hat{\mathbf{r}} \hat{R} \right) |\mathbf{r}\rangle = \hat{R}(\boldsymbol{\theta}) (\mathcal{R} \mathbf{r}) |\mathbf{r}\rangle.
\end{equation}
Since $|\mathbf{r}\rangle$ is an eigenstate of $\hat{\mathbf{r}}$ with eigenvalue $\mathbf{r}$, the rotated ket must be an eigenstate of the position operator with eigenvalue $\mathcal{R}\mathbf{r}$. Thus[cite: 55]:
\begin{equation}
    \hat{R}(\boldsymbol{\theta}) |\mathbf{r}\rangle = |\mathcal{R}(\boldsymbol{\theta})\mathbf{r}\rangle.
\end{equation}
Physically, this implies that the operator actively transports a particle localized at $\mathbf{r}$ to the rotated point $\mathcal{R}\mathbf{r}$.

\subsection{Infinitesimal Rotation}

\subsubsection{First-Order Expansion}
Consider an infinitesimal rotation by an angle $d\boldsymbol{\theta}$. Expanding the exponential to first order in the parameter $d\boldsymbol{\theta}$[cite: 59]:
\begin{equation}
    \hat{R}(d\boldsymbol{\theta}) \approx \hat{\mathbb{I}} - \frac{i}{\hbar} d\boldsymbol{\theta} \cdot \hat{\mathbf{L}}.
\end{equation}

\subsubsection{Action on the Wavefunction}
We examine how this operator transforms the wavefunction $\psi(\mathbf{r}) = \langle \mathbf{r} | \psi \rangle$. The transformed state $|\psi'\rangle = \hat{R}(d\boldsymbol{\theta})|\psi\rangle$ has the wavefunction:
\begin{equation}
    \psi'(\mathbf{r}) = \langle \mathbf{r} | \hat{R}(d\boldsymbol{\theta}) | \psi \rangle \approx \langle \mathbf{r} | \left( 1 - \frac{i}{\hbar} d\boldsymbol{\theta} \cdot \hat{\mathbf{L}} \right) | \psi \rangle.
\end{equation}
Distributing the bra-ket:
\begin{equation}
    (\hat{R}(d\boldsymbol{\theta})\psi)(\mathbf{r}) \approx \psi(\mathbf{r}) - \frac{i}{\hbar} d\boldsymbol{\theta} \cdot \langle \mathbf{r} | \hat{\mathbf{L}} | \psi \rangle.
\end{equation}
Using the position representation of angular momentum derived in Section 3, $(\hat{\mathbf{L}}\psi)(\mathbf{r}) = -i\hbar (\mathbf{r} \times \nabla)\psi(\mathbf{r})$, we substitute:
\begin{equation}
    (\hat{R}(d\boldsymbol{\theta})\psi)(\mathbf{r}) \approx \psi(\mathbf{r}) - \frac{i}{\hbar} d\boldsymbol{\theta} \cdot \left[ -i\hbar (\mathbf{r} \times \nabla) \psi(\mathbf{r}) \right].
\end{equation}
\begin{equation}
    (\hat{R}(d\boldsymbol{\theta})\psi)(\mathbf{r}) \approx \psi(\mathbf{r}) - d\boldsymbol{\theta} \cdot (\mathbf{r} \times \nabla) \psi(\mathbf{r}).
\end{equation}
Using the vector identity $\mathbf{A} \cdot (\mathbf{B} \times \mathbf{C}) = (\mathbf{A} \times \mathbf{B}) \cdot \mathbf{C}$, we rewrite the correction term:
\begin{equation}
    d\boldsymbol{\theta} \cdot (\mathbf{r} \times \nabla) = (d\boldsymbol{\theta} \times \mathbf{r}) \cdot \nabla.
\end{equation}
Recognizing that $d\mathbf{r} = d\boldsymbol{\theta} \times \mathbf{r}$ is the displacement vector due to rotation, we obtain the Taylor expansion for a scalar field under infinitesimal coordinate displacement[cite: 61]:
\begin{equation}
    \psi'(\mathbf{r}) \approx \psi(\mathbf{r}) - d\mathbf{r} \cdot \nabla \psi(\mathbf{r}) \approx \psi(\mathbf{r} - d\mathbf{r}).
\end{equation}
This confirms consistency with the passive transformation view: rotating the state active forward is equivalent to evaluating the original function at the reverse-rotated coordinates.

\subsection{Symmetries and Conservation Laws}

\subsubsection{Commutation with the Free Hamiltonian}
The free particle Hamiltonian is $\hat{H} = \frac{\hat{p}^2}{2m}$. To test for rotational invariance, we compute the commutator $[\hat{H}, \hat{L}_i]$[cite: 64].
Since $\hat{p}^2$ is a scalar operator (dot product of a vector operator with itself), it should commute with the generator of rotations. Explicitly:
\begin{equation}
    [\hat{L}_i, \hat{p}^2] = \sum_j [\hat{L}_i, \hat{p}_j \hat{p}_j] = \sum_j \left( \hat{p}_j [\hat{L}_i, \hat{p}_j] + [\hat{L}_i, \hat{p}_j] \hat{p}_j \right).
\end{equation}
Using the canonical commutation relation for angular momentum and vector operators, $[\hat{L}_i, \hat{p}_j] = i\hbar \epsilon_{ijk} \hat{p}_k$:
\begin{equation}
    [\hat{L}_i, \hat{p}^2] = i\hbar \sum_{j,k} \epsilon_{ijk} (\hat{p}_j \hat{p}_k + \hat{p}_k \hat{p}_j) = 2i\hbar \sum_{j,k} \epsilon_{ijk} \hat{p}_j \hat{p}_k.
\end{equation}
This sum involves the contraction of the antisymmetric tensor $\epsilon_{ijk}$ (in indices $j,k$) with the symmetric tensor $\hat{p}_j \hat{p}_k$. Therefore, the sum vanishes identically:
\begin{equation}
    [\hat{H}, \hat{\mathbf{L}}] = 0.
\end{equation}

\subsubsection{Conservation of Angular Momentum}
We invoke Ehrenfest's Theorem, established in Section 1[cite: 32]:
\begin{equation}
    \frac{d}{dt}\langle \hat{\mathbf{L}} \rangle = \frac{1}{i\hbar} \langle [\hat{\mathbf{L}}, \hat{H}] \rangle.
\end{equation}
Substituting the zero commutator derived above:
\begin{equation}
    \frac{d}{dt}\langle \hat{\mathbf{L}} \rangle = 0.
\end{equation}
This result [cite: 66] constitutes the quantum mechanical proof of the conservation of angular momentum for a free particle.

\subsubsection{Connection to the Correspondence Principle}
Classically, for a free particle, there is no torque ($\boldsymbol{\tau} = \mathbf{r} \times \mathbf{F} = 0$), implying $d\mathbf{L}_{cl}/dt = 0$.
The quantum result perfectly mirrors the classical law, as guaranteed by the Correspondence Principle. The rotational symmetry of space (isotropy) leads to the conservation of angular momentum in both frameworks, with the operator $\hat{\mathbf{L}}$ playing the dual role of the observable quantity and the generator of the symmetry group $SO(3)$.

\bibliography{./references.bib}
\nocite{*}

\end{document}
